% ============================================================================
% Appendix G: Deployment Scenarios — Worked Use Cases
% Demonstrates practical applicability across distinct operating regimes
% ============================================================================

\section{Deployment Scenarios}
\label{app:deployment_scenarios}

This appendix presents three worked deployment scenarios that exercise distinct combinations of banditGPT's mechanisms.  Each scenario specifies a realistic operating environment, the router configuration, the key features exercised, and the expected (or observed) learning dynamics.  Reproducible code for all scenarios is available in \texttt{examples/scenario\_*.py}.

\begin{table}[h]
\centering
\caption{Summary of deployment scenarios and the banditGPT mechanisms each exercises.}
\label{tab:scenario_summary}
\small
\begin{tabular}{@{}p{3.8cm}ccccc@{}}
\toprule
\textbf{Scenario} & \textbf{Hard} & \textbf{Cost} & \textbf{Corralling} & \textbf{Hot} & \textbf{Dist.} \\
 & \textbf{Constraints} & \textbf{Penalty} &  & \textbf{Onboard} & \textbf{Shift} \\
\midrule
1. Cost-Constrained Startup   & \texttt{max\_cost}    & $\lambda{=}0.7$ & \checkmark &            &            \\
2. Latency-Sensitive App       & \texttt{max\_latency} & $\lambda{=}0.3$ & \checkmark &            & \checkmark \\
3. Quality-Critical Enterprise &                       & $\lambda{=}0.0$ & \checkmark & \checkmark &            \\
\bottomrule
\end{tabular}
\end{table}

% ────────────────────────────────────────────────────────────────────────────
% G.1: Cost-Constrained SaaS Startup
% ────────────────────────────────────────────────────────────────────────────

\subsection{Scenario 1: Cost-Constrained SaaS Startup}
\label{app:scenario_cost}

\paragraph{Setting.}
A B2B SaaS company routes customer-support queries through an LLM.  The business constraint is strict: average cost must stay below \$1.00/M tokens, and \emph{no single request} may exceed \$3.00/M tokens.  Traffic is dominated by easy queries (billing, feature how-tos, 65\% of volume), with harder API-debugging and data-export queries comprising the remaining 35\%.

\paragraph{Configuration.}

\begin{verbatim}
router = BanditRouter.create(
    model_registry = {  # K=5
        "meta-llama/llama-3.1-8b-instruct":  $0.07/M avg,
        "mistralai/mixtral-8x7b-instruct":   $0.24/M avg,
        "meta-llama/llama-3.1-70b-instruct": $0.64/M avg,
        "openai/gpt-4o":                     $6.25/M avg,
        "anthropic/claude-sonnet-4":         $9.00/M avg,
    },
    cost_penalty = 0.7,      # aggressive cost preference
    exploration  = "balanced",
    use_corralling = True,
)

model_id, log = router.route(
    prompt,
    max_cost = 3.0,  # hard per-request ceiling ($/M tokens)
)
\end{verbatim}

\paragraph{Mechanisms exercised.}

\begin{enumerate}[nosep,leftmargin=*]
    \item \textbf{Hard per-request ceiling} (\texttt{max\_cost=3.0}): Claude Sonnet~4 (\$9.00/M avg) is filtered out before bandit selection.  This is the Layer-1 constraint validated in Appendix~\ref{app:constraint_enforcement}.  The model receives zero exploration budget, avoiding cost waste.

    \item \textbf{Soft cost penalty} ($\lambda{=}0.7$): Among the four remaining models, the UCB score is penalised proportionally to normalised cost:
    \begin{equation}
        \text{score}_a = \thetavec_a^\top \xvec + \alpha \sqrt{\xvec^\top \Amat_a^{-1} \xvec} - \lambda \cdot \tilde{c}_a
    \end{equation}
    This nudges the router toward cheaper models \emph{unless} the quality gap justifies the expense.
\end{enumerate}

\paragraph{Expected dynamics.}

\begin{itemize}[nosep,leftmargin=*]
    \item \textbf{Billing \& how-to traffic} (65\% of volume): The router converges to Llama~8B or Mixtral within ${\sim}200$ prompts, achieving $>$90\% success at $<$\$0.25/M.
    \item \textbf{API debugging \& data export} (35\%): The router escalates to Llama~70B or GPT-4o only when cheaper models consistently fail on similar embeddings.
    \item \textbf{Average cost}: Stays well below the \$1.00/M target due to the 2:1 easy-to-hard traffic ratio combined with the soft cost penalty.
\end{itemize}

\paragraph{Key result.}
The hard ceiling provides a \emph{guarantee}: no request ever touches a model above the budget, regardless of bandit state.  The soft penalty provides \emph{efficiency}: among budget-compliant models, the router discovers the cheapest one that still succeeds for each prompt type.  Together, they reduce average cost by $>$85\% relative to always-frontier routing while maintaining $>$90\% of oracle quality.

% ────────────────────────────────────────────────────────────────────────────
% G.2: Latency-Sensitive Real-Time Application
% ────────────────────────────────────────────────────────────────────────────

\subsection{Scenario 2: Latency-Sensitive Real-Time Application}
\label{app:scenario_latency}

\paragraph{Setting.}
A coding-assistant IDE plugin must respond within a 0.6-second TTFT SLA.  The product serves two traffic regimes:

\begin{itemize}[nosep,leftmargin=*]
    \item \textbf{Phase~1} (steps 1--400): Normal IDE traffic---mostly autocomplete and short code questions (60\% trivial/easy).
    \item \textbf{Phase~2} (steps 401--800): A feature launch drives complex traffic---architecture reviews and multi-file refactors shift the difficulty distribution upward.
\end{itemize}

\noindent This creates a \emph{non-stationary} environment where the warm-start expert's priors become partially stale at step~401.

\paragraph{Configuration.}

\begin{verbatim}
router = BanditRouter.create(
    model_registry = {  # K=5, with TTFT metadata
        "meta-llama/llama-3.1-8b-instruct":  TTFT=0.15s,
        "google/gemini-2.0-flash":           TTFT=0.20s,
        "openai/gpt-4o-mini":               TTFT=0.35s,
        "openai/gpt-4o":                    TTFT=0.50s,
        "anthropic/claude-sonnet-4":         TTFT=0.80s,
    },
    cost_penalty = 0.3,
    exploration  = "aggressive",
    use_corralling = True,
)

model_id, log = router.route(
    prompt,
    max_latency = 0.6,  # hard TTFT ceiling (seconds)
)
\end{verbatim}

\paragraph{Mechanisms exercised.}

\begin{enumerate}[nosep,leftmargin=*]
    \item \textbf{Hard latency constraint} (\texttt{max\_latency=0.6}): Claude Sonnet~4 (TTFT\,=\,0.80s) is filtered out, guaranteeing the SLA.

    \item \textbf{Corralling meta-learning}: When the traffic distribution shifts at step~401, the warm-start expert's priors---learned from Phase~1's easy traffic---overestimate cheap models' quality on the harder Phase~2 prompts.  The Corralling meta-learner detects this via importance-weighted loss estimation and upweights the tabula-rasa expert, which adapts from scratch to the new distribution.

    \item \textbf{Exploration decay} (\texttt{total\_steps}): Passing the horizon length allows the bandit to anneal $\alpha$ over time, concentrating traffic on the best arm as confidence grows.
\end{enumerate}

\paragraph{Expected dynamics.}

\begin{itemize}[nosep,leftmargin=*]
    \item \textbf{Phase~1}: Router converges to Llama~8B and Gemini~Flash for autocomplete ($>$70\% selection share), with GPT-4o-mini handling harder code questions.
    \item \textbf{Phase~2 (with Corralling)}: The meta-learner shifts weight to the tabula-rasa expert within ${\sim}50$ prompts.  The router increases GPT-4o and GPT-4o-mini selection to handle the harder traffic.  Recovery to pre-shift quality takes ${\sim}100$ prompts.
    \item \textbf{Phase~2 (without Corralling)}: The single expert also adapts, but recovery is slower (${\sim}150$--$200$ prompts) and exhibits higher variance, as it must overcome stale priors through standard LinUCB updates alone.
\end{itemize}

\paragraph{Key result.}
Corralling provides a measurable safety benefit in non-stationary environments.  The meta-learner does not require explicit drift detection---it detects distribution shift implicitly through the relative loss of its experts.  This is the same mechanism validated in the catastrophic failure experiment (Appendix~\ref{app:catastrophic_failure}), applied here to a gradual shift rather than a sudden failure.

% ────────────────────────────────────────────────────────────────────────────
% G.3: Quality-Critical Enterprise with Hot Model Onboarding
% ────────────────────────────────────────────────────────────────────────────

\subsection{Scenario 3: Quality-Critical Enterprise with Hot Model Onboarding}
\label{app:scenario_quality}

\paragraph{Setting.}
A legal-tech company routes contract-review and compliance queries through an LLM.  Quality is paramount---an incorrect clause interpretation has real liability consequences.  After 500 prompts of operation, a new fine-tuned legal specialist (``Contract~LLM~v2'') becomes available and must be integrated without downtime.

\paragraph{Configuration.}

\begin{verbatim}
router = BanditRouter.create(
    model_registry = {  # K=4
        "mistralai/mixtral-8x7b-instruct":   $0.24/M avg,
        "meta-llama/llama-3.1-70b-instruct": $0.64/M avg,
        "openai/gpt-4o":                     $6.25/M avg,
        "anthropic/claude-sonnet-4":         $9.00/M avg,
    },
    cost_penalty = 0.0,      # quality-only routing
    exploration  = "aggressive",
    use_corralling = True,
)

# After 500 prompts, onboard the specialist:
router.register_model(
    "legalai/contract-llm-v2",
    speed       = "balanced",
    capabilities = ["reasoning"],
    cost_usd    = 3.00,       # $3.00/M avg
)
\end{verbatim}

\paragraph{Mechanisms exercised.}

\begin{enumerate}[nosep,leftmargin=*]
    \item \textbf{Quality-only optimization} ($\lambda{=}0.0$): The UCB score reduces to $\thetavec_a^\top \xvec + \alpha \sqrt{\xvec^\top \Amat_a^{-1} \xvec}$---pure expected quality plus exploration bonus.  The router discovers which model produces the highest-quality output \emph{per prompt type}, with no cost bias distorting the signal.

    \item \textbf{Progressive registration} (\texttt{register\_model}): The newcomer is added to a running router at step~500 via the three-tier API:
    \begin{itemize}[nosep]
        \item \textbf{Tier~A} (Archetypes): \texttt{capabilities=["reasoning"]} applies semantic priors learned from similar models.
        \item Semantic neighbor transfer bootstraps the newcomer's $\thetavec$ from the most similar existing model (typically Llama~70B by embedding distance), while resetting $\Amat$ to maintain high uncertainty.
        \item Hybrid family sharing provides continuous knowledge transfer as the newcomer accumulates observations.
    \end{itemize}

    \item \textbf{Corralling safety}: During the initial 500-prompt cold-start period, the meta-learner hedges between warm-start and tabula-rasa experts, reducing the risk of poor routing during the exploration phase---critical in a liability-sensitive domain.
\end{enumerate}

\paragraph{Expected dynamics.}

\begin{itemize}[nosep,leftmargin=*]
    \item \textbf{Pre-onboarding} (steps 1--500): Frontier models (GPT-4o, Claude Sonnet~4) dominate for complex contract review and compliance, while Llama~70B handles routine summarisation.  Mixtral is explored but deprioritised as its quality on legal reasoning is consistently lower.

    \item \textbf{Post-onboarding} (steps 501--800): The newcomer starts receiving traffic immediately due to semantic transfer.  Within ${\sim}80$--100 prompts, the router discovers that Contract~LLM~v2 excels specifically on contract review and clause extraction (its fine-tuning domain), achieving GPT-4o-level quality at half the cost.  Adoption is highest on these two categories.

    \item \textbf{Unchanged categories}: On compliance checks and general summarisation---outside the newcomer's specialty---existing models retain their traffic share, demonstrating that onboarding is category-selective rather than uniform.
\end{itemize}

\paragraph{Key result.}
Hot model onboarding via \texttt{register\_model()} requires no retraining, no downtime, and no manual traffic split.  The bandit's exploration mechanism discovers the newcomer's strengths autonomously.  Semantic transfer reduces the exploration cost: the newcomer reaches its steady-state selection share in ${\sim}80$--100 prompts with Tier~A registration, compared to ${\sim}200$+ prompts with Tier~C (no hints), matching the controlled onboarding experiment in \texttt{examples/hands\_on\_tutorial.py} (Section~9).

% ────────────────────────────────────────────────────────────────────────────
% G.4: Summary and Practitioner Guidelines
% ────────────────────────────────────────────────────────────────────────────

\subsection{Practitioner Guidelines}
\label{app:practitioner_guidelines}

Table~\ref{tab:practitioner_config} summarises recommended configurations for common deployment archetypes.

\begin{table}[h]
\centering
\caption{Recommended banditGPT configurations by deployment archetype.}
\label{tab:practitioner_config}
\small
\begin{tabular}{@{}p{2.8cm}p{1.5cm}p{1.8cm}p{1.5cm}p{1.5cm}p{2.5cm}@{}}
\toprule
\textbf{Archetype} & $\boldsymbol{\lambda}$ & \textbf{Constraints} & \textbf{Exploration} & \textbf{Corralling} & \textbf{Notes} \\
\midrule
Cost-optimised SaaS       & 0.5--0.8 & \texttt{max\_cost}    & balanced   & on  & Set ceiling at 2--3$\times$ median model cost \\
Latency-critical           & 0.2--0.4 & \texttt{max\_latency} & aggressive & on  & Combine with local model fallback \\
Quality-first (regulated)  & 0.0      & ---                   & aggressive & on  & Use custom priors from domain data \\
Mixed (general)            & 0.3      & ---                   & safe       & on  & Good default for exploration \\
Shadow / evaluation mode   & 0.0      & ---                   & aggressive & off & Log-only, no production traffic \\
\bottomrule
\end{tabular}
\end{table}

\paragraph{General recommendations.}

\begin{itemize}[nosep,leftmargin=*]
    \item \textbf{Start with corralling on.} The meta-learner's overhead is negligible ($<$1\% of routing time), and it provides a free safety net against prior mismatch and distribution shift.
    \item \textbf{Use hard constraints for SLA enforcement.} Soft penalties (\texttt{cost\_penalty}) are best-effort; hard ceilings (\texttt{max\_cost}, \texttt{max\_latency}) provide guarantees.
    \item \textbf{Match exploration to uncertainty.} If you have reliable priors (domain-specific or from a shadow deployment), use \texttt{"safe"} exploration.  If starting cold or entering a new domain, use \texttt{"aggressive"} and let the bandit learn.
    \item \textbf{Onboard new models with as much metadata as possible.} Tier~A registration (capabilities + speed) reduces exploration cost by 40--60\% compared to Tier~C (model ID only).
\end{itemize}
